\documentclass[10pt]{article}
\providecommand{\FIGDIR}{fig}
\newcommand{\DIGESTDATE}{30 September 2026}
\input{preamble}

\begin{document}

% ==================================================================== MASTHEAD
\thispagestyle{empty}
\begin{tikzpicture}[remember picture, overlay]
  \fill[Deep] (current page.north west) rectangle
      ($(current page.north east)+(0,-67mm)$);
  \fill[Coral] ($(current page.north west)+(0,-67mm)$) rectangle
      ($(current page.north east)+(0,-68.4mm)$);
  \foreach \x/\y/\r in {22/12/0.9, 41/26/1.5, 63/9/0.7, 88/31/1.1, 112/17/1.9,
                        138/29/0.8, 159/13/1.3, 178/54/1.0, 196/21/1.6,
                        30/44/0.6, 74/57/1.2, 125/48/0.9, 168/6/0.8, 191/59/0.7,
                        52/18/1.0, 100/60/0.6, 148/20/1.1, 66/49/1.4, 205/45/0.8}
    \fill[Sky, opacity=0.30] ($(current page.north west)+(\x mm,-\y mm)$) circle (\r mm);
\end{tikzpicture}

\vspace*{4mm}
{\sffamily\fontsize{7.6}{9}\selectfont\color{Coral}%
 MONTHLY ROLLUP\;\textbullet\;PhD RESEARCH WATCH\par}
\vspace{2.2mm}
{\sffamily\bfseries\fontsize{27}{29}\selectfont\color{white}%
 Particulate Matter\par}
{\sffamily\fontsize{27}{29}\selectfont\color{Sky}%
 Monitoring \& Health Effects\par}
\vspace{5.0mm}
{\sffamily\fontsize{8.4}{10}\selectfont\color{white}%
 September 2026\;\;\textcolor{Coral}{\textbar}\;\;Entry window 1--30 September 2026, 30 daily issues\;\;%
 \textcolor{Coral}{\textbar}\;\;622 records pooled\par}
\vspace{18mm}

% -------------------------------------------------------------- metric strip
\noindent
\begin{minipage}[t]{0.190\linewidth}\metric{622}{RECORDS POOLED\\(30 DAILY ISSUES)}{Deep}\end{minipage}\hfill
\begin{minipage}[t]{0.190\linewidth}\metric{148}{EXPOSURE ASSESSMENT\\(24\%), LARGEST BAND}{Teal}\end{minipage}\hfill
\begin{minipage}[t]{0.190\linewidth}\metric{173}{EFFECT ESTIMATES\\FROM 93 SOURCES}{Sky}\end{minipage}\hfill
\begin{minipage}[t]{0.190\linewidth}\metric{143}{TIER-A\\RECORDS (23\%)}{Amber}\end{minipage}\hfill
\begin{minipage}[t]{0.190\linewidth}\metric{98}{METADATA-ONLY\\RECORDS (16\%)}{Clay}\end{minipage}

\vspace{6mm}

% ------------------------------------------------------------ signal of month
\section*{\sffamily\bfseries\color{Deep}Signal of the month}
\vspace{-1.5mm}{\color{Coral}\rule{\linewidth}{1.1pt}}\vspace{2.5mm}

\begin{keybox}
{\sffamily\bfseries\small\color{Deep}September settled one question at population scale: per
unit mass, fire-sourced PM$_{2.5}$ is more lethal than background PM$_{2.5}$, and the excess sits
in its carbonaceous fraction --- the part a mass-reporting optical network does not resolve.}\\[1.4mm]
\footnotesize
Three large, methodologically independent studies landed inside sixteen days. Su et al.\
(\textit{Lancet Oncol}, 15 Sep) followed 414,016 SEER--Medicare lung-cancer patients: per
1\,$\mu$g\,m$^{-3}$, death HR \textbf{1.078} (1.064--1.092) for wildfire against \textbf{1.017}
(1.015--1.019) for non-wildfire PM$_{2.5}$. Ju et al.\ (\textit{Lancet Planet Health}, 16 Sep), with
instrumental-variable models across 1,366 communities in ten countries, found landscape-fire
particles are \textbf{19.0\,\%} of ambient PM$_{2.5}$ mass but \textbf{30.6\,\%} of the attributable deaths,
and that conventional time-series estimates understate the fire effect by more than half. Jin et
al.\ (\textit{ES\&T}, 30 Sep) split smoke PM$_{2.5}$ by composition in 77.5 million Medicare
beneficiaries: HR \textbf{1.032} per 1\,$\mu$g\,m$^{-3}$ for smoke mass and \textbf{1.087} for smoke carbon,
with the carbon fraction carrying most of the 22,420 attributable deaths a year.
\textbf{The common caveat is arithmetic}: chronic smoke exposure is small and narrow (median
0.28\,$\mu$g\,m$^{-3}$ in Su et al.), so ``per 1\,$\mu$g\,m$^{-3}$'' coefficients extrapolate across most
of their own range, and on the IQR scale the ranking can reverse. The direction, though, now holds
across a cancer cohort, a multi-country IV design and a composition-resolved national cohort.
Kaur et al.'s twenty-year western-US source apportionment (\textit{Atmos Pollut Res}, 20 Sep) adds
the measurement side: the mass that is disappearing is secondary sulfate and nitrate, which
optical sensors see well, while the share from residential-wood and fire combustion is rising.
\textbf{The month's verdict: for health-relevant exposure in fire-affected regions, a network
that adds an absorption or BC channel buys more than one that adds sensor density.}
\end{keybox}

\vspace{3mm}

\section*{\sffamily\bfseries\color{Deep}What changed in September}
\vspace{-1.5mm}{\color{Coral}\rule{\linewidth}{1.1pt}}\vspace{2.5mm}

\footnotesize
\begin{itemize}
\item \textbf{Calibration moved from ``does it fit'' to ``is the fit honest''.} Sahin et al.\
(\textit{Atmosphere}, 29 Sep) showed six ML calibrators reaching R$^2$ 0.86--0.89 under random
splits and near or below zero under chronological hold-out, with validation design explaining
\textbf{63--91\,\%} of the R$^2$ spread and algorithm choice 3.6--10.6\,\%. Kaur, Mangin \& Kelly
(\textit{Aerosol Sci Technol}, 3 Sep) showed Plantower's PMSX003N fixes dust (PM$_{2.5}$ slope
1.85--1.91 vs 0.15 for the PMS5003T) by re-binning $>$1\,$\mu$m particles, which is the same mechanism
that makes it over-read in humid air. Khalili Hollo et al.\ (\textit{ACS Environ Au}, 25 Sep) put
fifty Atmotubes and five QuantAQ units through four co-locations in two years: unit response
\textbf{drifted between co-locations}, 5--11\,\% of wearables were dead at any comparison, and RH
correction cut RMSE 42--50\,\%. Together these define the operating rule: temporally blocked
validation, re-co-location rather than a single pre-deployment fit, and slope reported with RMSE.

\item \textbf{Neither side of the co-location is stationary.} Zhou et al.\ (\textit{PNAS}, 14 Sep)
used the 2025 suspension of the US embassy monitoring programme as a natural experiment: local
satellite AOD rose \textbf{24--33\,\%} after the monitors went dark, most where no substitute data
existed --- monitoring changed behaviour, not just the record. A Korean analysis (25 Sep) is
titled on step changes in the regulatory PM record when stations relocate, and a mobile monitor
250\,m from a regulatory station published a different AQI colour one time in ten (7 Sep).

\item \textbf{August's prediction (1) was met.} A field evaluation of three community-priced
condensation particle counters at an urban regulatory station (\textit{Aerosol Sci Technol},
12 Sep) gave \textbf{R$^2>0.99$} against a research-grade CPC with a 3\,\% inter-unit CV over four
months. The Partector 2 Pro matched SMPS/CPC number concentration within $\pm$10\,\% but under-sized
by up to 28\,\% (\textit{J Occup Environ Hyg}, 15 Sep), and a six-channel electro-inertial sizer
reproduced SMPS distributions within 10\,\% over a 16-day rooftop deployment (Le et al.,
\textit{Meas Sci Technol}, 22 Sep). None ran longer than one season, so drift is still unmeasured.
\textbf{Prediction (2) --- an ageing term inside an operational network correction --- was not
met}; the Denver campaign corrects drift by re-co-location, not by a modelled ageing term.

\item \textbf{Epidemiological method got sharper.} Leung, Neophytou \& Wilson (\textit{Int J
Epidemiol}, 13 Sep) proved that zero-filling gestational exposure histories in distributed-lag
models \emph{manufactures} the late-pregnancy window it is used to detect. A 115,056-birth
analysis returned an informative null for short-term PM$_{2.5}$ and spontaneous preterm birth (17
Sep). In the Sichuan Basin, excess cardiorespiratory mortality per 10\,$\mu$g\,m$^{-3}$ rose from 0.4\,\%
to \textbf{1.4\,\%} as concentrations fell, handing back about \textbf{19\,\%} of the Clean Air Action's
avoided deaths (6 Sep) --- the supralinear low-concentration slope showing up as a policy number.

\item \textbf{Personal and continuous metrics carried signals area metrics missed.} Time above
35\,$\mu$g\,m$^{-3}$ indoors, not the period mean, carried the cystic-fibrosis pilot signal (26
Sep); a personal-monitor panel in older adults found a next-day activity deficit of about 8 min per
IQR that area PM did not reproduce (27 Sep); and a third of indoor PM$_{2.5}$ variance sat between
homes after full adjustment (8 Sep). Against that, \textbf{two African maternal cohorts} (21--22
Sep) rested exposure on uncorrected consumer optical sensors in biomass smoke --- the cheapest
sensor in the hardest aerosol.

\item \textbf{Composition outran mass across the toxicology band.} Brake dust as a catalyst
(10 Sep), saturating rather than additive oxidative potential (19 Sep), \textit{p}-phenylenediamine
antioxidants in open waste burning at 11.5\,$\mu$g\,g$^{-1}$ PM$_{2.5}$ (29 Sep) and tyre-related
chemicals in Leipzig roadside PM$_{10}$ (30 Sep). None of these is visible to a mass sensor.

\item \textbf{Occupational \& indoor nearly doubled} (37 $\rightarrow$ 72 records) and
\textbf{Cardiovascular \& metabolic doubled} (23 $\rightarrow$ 46), while \textbf{Sensing fell} from
140 to 116 (25\,\% $\rightarrow$ 19\,\%) as the AMT/ACP Crossref channel delivered 32 records against
48 in August. Exposure assessment (148, 24\,\%) is now the largest band.
\end{itemize}

\clearpage

% ==================================================================== FIGURES
\section*{\sffamily\bfseries\color{Deep}Corpus analytics --- September 2026}
\vspace{-1.5mm}{\color{Coral}\rule{\linewidth}{1.1pt}}\vspace{2.5mm}

\noindent\includegraphics[width=\linewidth]{\FIGDIR/f1_subtopics.png}
\figcap{Subtopic distribution across all 622 pooled records; all ten subtopics populated.
\textit{Exposure assessment \& modelling} 148 (24\,\%) and \textit{Sensing, forecasting \&
instrumentation} 116 (19\,\%) together hold 43\,\%, down from August's 48\,\%. \textit{Occupational \&
indoor} 72, \textit{Mechanistic toxicology} and \textit{Burden, policy \& mitigation} 55 each,
\textit{Cardiovascular \& metabolic} 46, \textit{Respiratory \& allergic} 44, \textit{Neuro / mental
health} 30, \textit{Reproductive \& developmental} 29, \textit{Other clinical} 27. Single-label
assignment by dominant contribution.}

\vspace{3mm}

\noindent\includegraphics[width=\linewidth]{\FIGDIR/w1_subtopic_trend.png}
\figcap{Records per Sunday--Saturday week by subtopic across the 30 September issues (the
first and last weeks are partial): 120 for 1--5 Sep, then 126, 148 and 162, and 66 for 27--30 Sep.
Daily volume ranged from 4 (27 Sep, a Sunday) to 38 (3 Sep) and tracks publisher deposit timing, so
week totals are the readable unit. Volume rose through the month; the 26--28 Sep issues were built
late as a backfill but each kept its own entry-date window. This panel replaces the July-era
``what the daily cadence missed'' chart, which had no meaning once every day had its own issue and
was shipped with a stale July legend in the August monthly.}

\vspace{3mm}
\noindent\includegraphics[width=\linewidth]{\FIGDIR/f2a_architecture.png}
\figcap{Study architecture. Measurement campaigns 118, metadata-only 98, modelling/inventory 88,
reviews 65, cohorts 61, experimental/toxicology 59, cross-sectional 48, acute designs 33,
chamber/laboratory 22, ecological 18. \textbf{Metadata-only rose from 12\,\% in August to 16\,\%},
almost all Elsevier deposits without abstracts, while Scholar Gateway (the full-text route)
has been disconnected since 29 Sep. \textbf{Trials or interventions: 8 records in 622} --- double
August's four, still about 1\,\%.}

\clearpage

%% f2b_endpoint.png is DELIBERATELY OMITTED again (2026-10-01). ENDPOINT_CANON now exists,
%% but at month scale it still returns 48 categories; the figure is 3958x5843 px and would
%% render ~48 rows at an illegible size on one page. Distribution given as text.
\begin{keybox}
{\sffamily\bfseries\footnotesize\color{Deep}Health endpoint distribution --- reported as
text, not as a figure}\\[1.2mm]
\footnotesize
\texttt{ENDPOINT\_CANON} now exists (August's 144 raw strings collapse to a bounded map), but at
month scale it still yields \textbf{48 categories}, so the pooled chart would be 48 rows at an
unreadable size. \textbf{303 records (49\,\%) carry no human health endpoint} --- instrument,
exposure, emissions and atmospheric-process work plus the metadata-only records. Of the
remainder: \textbf{respiratory 66, cardiovascular 33, other clinical 23, neuro 28} (cognitive 17,
mental health 11), \textbf{cancer 26} (still split as ``Oncologic'' 15 and ``Cancer'' 11 --- the open
canonicalisation defect), reproductive 17, mortality 18, metabolic 11, in-vitro toxicological 12,
neurodevelopmental 6, cerebrovascular 6, with the balance in one- to five-record categories.
\textbf{Fix still owed:} merge Cancer/Oncologic and the two neuro labels, and cap the map at
$\sim$20 categories so f2b can return at rollup scale.
\end{keybox}

\vspace{3mm}

\noindent\includegraphics[width=\linewidth]{\FIGDIR/f3_metric_geography.png}
\figcap{Left --- particle metric. Unspeciated PM or emissions 243, PM$_{2.5}$ alone 202, PM$_{2.5}$ with
PM$_{10}$ 84, optical properties 30, composition/speciation 19, PM$_{10}$ alone 13, \textbf{ultrafine
10 and size-resolved number 8} (3\,\% together, against 4\,\% in August), composite indices 8,
bioaerosol 5. Right --- geography. Global / multi-region 190 (inflated by metadata-only records with
no recoverable site), \textbf{China 165 (27\,\%)}, Europe 76, North America 68, East Asia ex-China 38,
South Asia 25, \textbf{Sub-Saharan Africa 22} (16 in August), Middle East \& N.\ Africa 13, Southeast
Asia 9, Latin America 9 (13 in August), Oceania 4, polar/remote 2, Central Asia 1.}

\clearpage

\noindent\includegraphics[width=\linewidth]{\FIGDIR/f4_heatmap.png}
\figcap{Subtopic $\times$ study-architecture cross-tabulation. The dense cells are the expected
ones --- sensing $\times$ measurement campaign, exposure $\times$ modelling and metadata-only. The
informative absences are unchanged: the trial column is nearly empty in every health subtopic, and
the sensing and exposure rows contribute almost nothing to the effect-estimate pool.}

\vspace{3mm}

\noindent\includegraphics[width=\linewidth]{\FIGDIR/f6_lifecourse.png}
\figcap{Life-course windows across the month (non-exclusive; only records with an identifiable
human window): in utero 41, infancy/childhood 35, adolescence 10, working age 77, older adults 40.
\textbf{Adolescence remains the thinnest window}, as in July and August.}

\clearpage

\noindent\includegraphics[width=\linewidth,height=0.80\textheight,keepaspectratio]{\FIGDIR/f5_forest.png}
\figcap{Quantitative effect estimates extracted across September, panel 1 of 6. \textbf{This is not a
meta-analysis and no pooled summary is drawn.} Contrasts are not commensurable (per 1, per 10, per
IQR, per doubling, threshold, quartile, fuel-type and occupational contrasts) and the outcome scale
mixes OR, HR, RR, IRR, PR, sHR and ratio-of-means across unrelated diseases, so
\textbf{heterogeneity is not estimable and none is reported}. Of 173 estimates, 161 carry a CI on a
ratio or difference scale: \textbf{132 exclude the null in the harmful direction, 10 in the
protective direction and 19 cross it.}}

\clearpage
\noindent\includegraphics[width=\linewidth,height=0.86\textheight,keepaspectratio]{\FIGDIR/f5_forest_2.png}
\figcap{Effect estimates, panel 2 of 6. Same caveats as panel 1.}
\clearpage
\noindent\includegraphics[width=\linewidth,height=0.86\textheight,keepaspectratio]{\FIGDIR/f5_forest_3.png}
\figcap{Effect estimates, panel 3 of 6. Same caveats as panel 1.}
\clearpage
\noindent\includegraphics[width=\linewidth,height=0.86\textheight,keepaspectratio]{\FIGDIR/f5_forest_4.png}
\figcap{Effect estimates, panel 4 of 6. Same caveats as panel 1.}
\clearpage
\noindent\includegraphics[width=\linewidth,height=0.86\textheight,keepaspectratio]{\FIGDIR/f5_forest_5.png}
\figcap{Effect estimates, panel 5 of 6. Same caveats as panel 1.}
\clearpage
\noindent\includegraphics[width=\linewidth,height=0.80\textheight,keepaspectratio]{\FIGDIR/f5_forest_6.png}
\figcap{Effect estimates, panel 6 of 6. None of the ten estimates below the null is a protective
PM effect: four are protective \emph{co-exposures} (greenness against VTE, green space and natural
environment against COPD, exercise against mortality) and six are harms on a decreasing scale
(peak-flow ratio, femur length, abdominal circumference, two next-day activity betas, and
scholarly output per $\mu$g\,m$^{-3}$). The largest
ratios (OR 4.78 for type A aortic dissection per doubling of same-day PM$_{2.5}$; OR 4.09 for
microplastic detection in lung malignancy) come from small case series with wide intervals.}

\clearpage

\section*{\sffamily\bfseries\color{Deep}Emerging, fading, and recurring}
\vspace{-1.5mm}{\color{Coral}\rule{\linewidth}{1.1pt}}\vspace{2.5mm}

\footnotesize
\textbf{Emerging.} \textbf{(i) Source-specific potency as the default framing for fire PM} ---
see the signal box; three large designs agree on direction. \textbf{(ii) Validation honesty in
sensor calibration} --- temporally blocked hold-outs (Sahin et al.), stratified error reporting
across AQI bins (the Vishakhapatnam co-location, 15 Sep), and repeated co-location to track drift
(Khalili Hollo et al.). \textbf{(iii) Community-cost particle-number instruments} reaching
reference agreement (cCPC, Partector 2 Pro, electro-inertial sizer). \textbf{(iv) Monitoring as an
intervention} --- the embassy-shutdown natural experiment is the first record in this corpus that
treats a monitor's existence as the exposure. \textbf{(v) Continuous-sensor-only metrics}
(time-above-threshold, next-day personal exposure) carrying signals that means and area estimates
miss.

\vspace{2mm}
\textbf{Fading.} Two complete months are still too few to call anything fading, and no decline
is asserted. The observation carried from August stands: \textbf{single-pollutant time-series
studies} continue to report PM as a weak term once co-pollutants enter the model --- the Japanese
cancer-mortality series (30 Sep) lost significance after NO$_2$ or SO$_2$ adjustment. The share of
\textit{Sensing} fell from 25\,\% to 19\,\%, but that tracks the Crossref AMT/ACP channel (48
$\rightarrow$ 32 records), not the literature.

\vspace{2mm}
\textbf{Recurring groups and venues.} \textbf{No author-cluster analysis is possible} --- the corpus
has no affiliation field (open since 15 August), so surname repeats in a China-heavy corpus are not
investigator clusters. Venue concentration is supportable: \textit{Environ Sci Technol} 46,
\textit{Environ Pollut} 36, \textit{Front Public Health} 28, \textit{Atmosphere} 26, \textit{Atmos
Environ} 24, \textit{Atmos Chem Phys} 23, \textit{Atmos Pollut Res} 15, \textit{Environ Res} 14,
\textit{J Hazard Mater} 13, \textit{Indoor Air} 12, \textit{Atmos Meas Tech} 9. \textit{ES\&T}
displaced \textit{ACP} as the largest venue. \textbf{Most-cited work of the period is not reported}:
citation counts on records days to weeks old measure indexing lag, not influence. The month's
highest-venue records were Su et al.\ (\textit{Lancet Oncol}), Ju et al.\ (\textit{Lancet Planet
Health}) and Zhou et al.\ (\textit{PNAS}).

\clearpage

% ------------------------------------------------------------------ trial watch
\section*{\sffamily\bfseries\color{Deep}Trial registry watch}
\vspace{-1.5mm}{\color{Coral}\rule{\linewidth}{1.1pt}}\vspace{2.5mm}

\footnotesize
Daily registry sweeps (\texttt{LastUpdatePostDate} windows) ran through September; \texttt{state/trials.json} holds
\textbf{16 trials}, \textbf{seven} of them first seen in September (the August rollup reported 24; the
difference is not reconciled in the run log --- see defect (vi)). The recurring finding of August hardened into a pattern: \textbf{registrations framed
around air pollution that register no exposure measurement.}

\vspace{1.5mm}
\textbf{NCT07817836} (IDEA, Pavia, observational, $n$=151): urban outdoor workers vs office controls,
personal passive sampling, but no PM$_{2.5}$ outcome measure. \textbf{NCT07713264} (VISIONS, Stanford,
interventional, $n$=50, outdoor workers): \texttt{analyze\_endpoints} shows three psychosocial scales
and no air-quality measure. \textbf{NCT07793123} (interventional, $n$=60, placebo-controlled): the
third consecutive topical hit whose PM framing no measured variable carries. \textbf{NCT05423665}
(observational, $n$=360) lists air pollution as a condition beside fetal growth restriction.
\textbf{NCT07196436} (AIRVAC, energy-recovery ventilation plus air cleaner, $n$=80) was found through
its protocol paper rather than the sweep; registry co-primaries are spirometry and asthma control,
the protocol primary is asthma control alone. \textbf{NCT06541691} (COATED-AIR, $n$=3,000) has no
exposure endpoint and a stale recruiting status. \textbf{NCT02944682} (HAPIN, $n$=6,704, LPG
intervention) posted a registry update; main results are published.

\vspace{1.5mm}
Published trial-class records rose to \textbf{8} in 622, the most instrument-relevant being the
protocol for a Dutch primary-school trial (NCT07479420) randomising \textbf{180 classrooms in 29
schools} to HEPA, ionisation/plasma or no cleaner (9 Sep) and a smoke-free-home RCT. The 30 Sep
sweep returned one observational COPD-exacerbation cohort (CAN-EXACT, NCT07850648), screened out:
environmental factors appear only inside a generic secondary measure. \textbf{The interventional
evidence base for personal and household PM reduction is still thickening in the registry faster
than in the literature, and the registry entries increasingly omit the exposure they claim to
target.}

\vspace{4mm}

% -------------------------------------------------------------------- open gaps
\section*{\sffamily\bfseries\color{Deep}Open gaps and read on direction}
\vspace{-1.5mm}{\color{Coral}\rule{\linewidth}{1.1pt}}\vspace{2.5mm}

\begin{keybox}
{\sffamily\bfseries\footnotesize\color{Deep}Gaps in the literature}\\[1.2mm]
\footnotesize
\textbf{(i) Composition is where the health signal is, and networks report mass.} Fire-carbon
potency, oxidative-potential saturation and quinone-forming antioxidants all argue for absorption,
BC or OP channels; no September record deployed one at network scale with a field calibration.
\textbf{(ii) Drift is measured but not modelled.} Denver shows drift between co-locations; no record
places an ageing term inside an operational correction (August's prediction 2, re-issued).
\textbf{(iii) Ultrafine instruments now pass co-location but have no long-term record}: every
September PN validation ran for a season or less.
\textbf{(iv) Low-cost sensors in biomass smoke without correction} underpin new maternal-health
cohorts in Sub-Saharan Africa --- the region whose share of the corpus is growing (22 records) and
whose aerosol is the worst case for optical sensing.
\textbf{(v) Trials without exposure measurement.} Eight trial-class records; registry entries
increasingly name air pollution without measuring it.
\end{keybox}

\vspace{2.5mm}

\begin{keybox}
{\sffamily\bfseries\footnotesize\color{Deep}Gaps and defects this pipeline creates}\\[1.2mm]
\footnotesize
\textbf{(i) No duplicates this month.} 622 records, \textbf{622 unique DOIs}; the
\texttt{DUPLICATE-DOI} hard-fail added on 1 September held for all 30 issues.
\textbf{(ii) 98 records (16\,\%) are metadata-only}, up from 12\,\%. They are counted but carry no
assessed content. The rise tracks Elsevier deposit practice and the loss of Scholar Gateway
(``could not resolve user identity'' since 29 Sep; needs reconnection).
\textbf{(iii) Endpoint canonicalisation is incomplete} (48 categories; Cancer/Oncologic split).
\textbf{(iv) No affiliation field} --- no author or institution clusters.
\textbf{(v) Source legs.} OpenAlex returned 429/403 all month; Europe PMC was down for the whole 25
Sep run (HTTP 503); arXiv returned 406 from mid-September and recovered on 29 Sep; Consensus
returned \textbf{0 new} records on every September sweep. Three daily runs (26--28 Sep) did not fire
and were backfilled day by day on 29 Sep, each with its own window.
\textbf{(vi) Trial-list continuity.} \texttt{trials.json} lists 16 trials against the 24 the August
rollup reported; earlier entries appear to have been dropped without a log line. To audit before
the October rollup.
\end{keybox}

\vspace{3mm}

\footnotesize
\textbf{Read on direction.} August's read was that the two halves of the field were optimising
different objective functions --- sensing on how far a PM$_{2.5}$ mass correction transfers,
epidemiology on toxicity and number. September narrows that gap from both sides. Epidemiology
supplied the strongest population-scale evidence yet that source and composition, not mass, set
potency; instrumentation supplied the first community-cost particle counters with reference
agreement and, more usefully, a set of papers that make calibration claims auditable (blocked
validation, re-co-location, stratified error). \textbf{What has not moved is deployment}: no network
in the corpus yet measures the fraction the epidemiology now points to.
\textbf{Falsifiable predictions for October:}
(1) a low-cost network with an absorption/BC or particle-number channel and a field calibration of
at least one full season appears; and
(2) at least one sensor-calibration paper reports a temporally blocked or prospective hold-out as
its primary metric, citing the validation-leakage argument. If (2) fails, random-split R$^2$ remains
the field's reporting norm and should be discounted by default.

\clearpage

% ==================================================================== DIGEST
%% MANDATORY, same contract as the daily and the weekly (formalised 2026-08-02).
%% Full summaries are lifted VERBATIM from issues/digest_*.tex by mkdigest.py.
\section*{\sffamily\bfseries\color{Deep}Paper-level digest --- the month by subtopic}
\vspace{-1.5mm}{\color{Coral}\rule{\linewidth}{1.1pt}}\vspace{2mm}

\footnotesize
Every pooled record appears below. Records that carry a full assessed summary have it lifted
verbatim from the daily issue that first published it; a title-only or extractive entry is
\textbf{not an assessment} and the absence of criticism is not approval. Metadata-only records are a
pointer to a DOI and nothing more.

\vspace{3mm}

\input{digest_body}

\clearpage

% ==================================================================== TABLE
\section*{\sffamily\bfseries\color{Deep}Pooled corpus register --- September 2026}
\vspace{-1.5mm}{\color{Coral}\rule{\linewidth}{1.1pt}}\vspace{2.5mm}

\input{table_rows}

\vspace{3mm}

\begin{keybox}
{\sffamily\bfseries\footnotesize\color{Deep}Provenance and method}\\[1.2mm]
{\fontsize{7.2}{8.8}\selectfont
Pure aggregation over \texttt{state/corpus/2026-09-*.json} (30 daily issues, entry windows 1--30
September 2026, contiguous), \texttt{state/metrics.csv} and \texttt{state/trials.json}. No API was
re-queried and no DOI re-resolved; every record passed the DOI gate at daily-issue time. Figures are
pooled re-renders of the daily set; the forest panels catalogue estimates and do not pool them.
Abstracts remain the copyright of their respective publishers.\par}
\end{keybox}

\end{document}
